% =====================================================================
% Ensaio 1 — Seção 2: Background
%
% Bellemare dispensa a seção de Background em artigo de journal, mas a considera
% útil "quando os detalhes de uma legislação precisam ser tidos em mente para
% avaliar os efeitos dessa legislação sobre algum desfecho". É exatamente o caso.
%
% ⚠️ CORREÇÃO SUBSTANTIVA em relação ao esboço do projeto: o framework é centrado
% em GLIFOSATO. A química documentada na Chapada do Apodi NÃO é essa. Ver §2.3.
% =====================================================================

\section{Background}
\label{sec:background}

\subsection{A Lei Zé Maria do Tomé e sua cronologia}
\label{sub:lei}

Em 08 de janeiro de 2019 o governador do Ceará sancionou a Lei Estadual nº
16.820, publicada no Diário Oficial e vigente a partir de 09 de janeiro do mesmo
ano.\footnote{A distinção entre sanção e vigência é irrelevante em painel mensal
--- ambas caem em janeiro de 2019 ---, mas é registrada aqui porque o texto
deve dizer 09/01/2019 ao falar de vigência e 08/01/2019 ao falar de sanção.} A lei
acrescenta o art. 28-B à Lei Estadual nº 12.228/1993, norma-mãe sobre agrotóxicos
no estado, com a seguinte redação:

\begin{quote}
\textbf{Art. 28-B.} É vedada a pulverização aérea de agrotóxicos na agricultura
no Estado do Ceará.

\textbf{§ 1º} A infração ao art. 1º sujeita o infrator ao pagamento de multa de
15 mil (quinze mil) UFIRCEs.

\textbf{§ 2º} Fica proibida a incorporação de mecanismos de controle vetorial por
meio de dispersão por aeronave em todo o Estado do Ceará, inclusive para os casos
de doenças causadas por vírus.
\end{quote}

Dois traços da redação organizam todo o desenho empírico. Primeiro, a proibição é
de \textbf{método}: nenhum princípio ativo é nomeado. Segundo, ela não se
restringe à agricultura --- o §2º alcança a dispersão aérea sanitária, o que
significa que municípios sem agricultura pulverizada, mas com controle aéreo de
endemias, também são alcançados pela norma.

O Ceará foi o primeiro estado brasileiro a adotar a medida. A lei leva o nome de
José Maria Filho, conhecido como Zé Maria do Tomé, líder comunitário da Chapada
do Apodi assassinado em 2010 após denunciar publicamente a pulverização aérea na
região. O dado não é ornamental: ele situa a proibição como resposta a um
conflito distributivo, e não como ajuste técnico --- ponto ao qual a
§\ref{sec:teoria} retorna.

\begin{table}[htbp]
\centering
\caption{Cronologia legal do tratamento}
\label{tab:cronologia}
\small
\begin{tabular}{lp{10.5cm}}
\hline
Data & Evento \\
\hline
09/12/1993 & Lei Estadual nº 12.228/1993 --- norma-mãe sobre agrotóxicos no Ceará. \\
2009 & Limoeiro do Norte proíbe a pulverização aérea por lei municipal, ao amparo do art. 29 da lei estadual, que autoriza legislação supletiva. \\
2010 & Assassinato de Zé Maria do Tomé. \\
\textbf{24/02/2015} & \textbf{Apresentação do Projeto de Lei nº 18/2015.} Primeiro sinal público --- \emph{marco de notícia}. \\
\textbf{18/12/2018} & \textbf{Aprovação por unanimidade na Assembleia Legislativa}, após quatro anos de tramitação --- \emph{marco de certeza}. \\
\textbf{09/01/2019} & \textbf{Vigência da Lei nº 16.820/2019} --- \emph{marco de obrigação}. \\
mai/2023 & O Supremo Tribunal Federal julga a lei constitucional (ADI 6137, unânime). \\
19/12/2024 & Lei Estadual nº 19.135/2024 dá nova redação ao art. 28-B, abrindo exceção para aeronaves remotamente pilotadas. \\
2025 & Nova contestação no STF (ADI 7794), pendente. \\
\hline
\end{tabular}
\end{table}

\paragraph{Três marcos, não um.} A Tabela~\ref{tab:cronologia} distingue
deliberadamente notícia, certeza e obrigação. A distinção tem consequência
empírica direta: a janela de medida da dose pré-ban começa \emph{depois} do
primeiro marco, de modo que a variável de tratamento pode já incorporar resposta
antecipatória. A §\ref{sub:dose} desenvolve o ponto.

\paragraph{A reversão de 2024, e por que ela fecha a janela.} A Lei nº
19.135/2024 mantém a vedação, \emph{"salvo se realizada por meio de Aeronaves
Remotamente Pilotadas --- ARPs, Veículo Aéreo Não Tripulado --- VANT ou Drones"},
sob condições de altura, velocidade de vento e distância de escolas, hospitais e
áreas protegidas. A nova redação \textbf{não reproduz} o §2º de 2019: a proibição
de dispersão aérea para controle vetorial cai. Do ponto de vista de instrumento
de política, o Ceará não revogou a proibição --- \textbf{substituiu o
instrumento}, trocando cota zero por regulação de desempenho com zona-tampão.
Para o Ensaio 1, o efeito é fixar o fim da janela de análise em 19/12/2024, de
modo a preservar a cota zero como tratamento.

\subsection{A Chapada do Apodi e a fruticultura irrigada}
\label{sub:chapada}

A Chapada do Apodi, no baixo Jaguaribe, concentra o agronegócio de fruticultura
irrigada do Ceará --- notadamente banana e melão, em perímetros de irrigação
sobre solo cárstico. É a região à qual a lei responde, e é onde a exposição
populacional à deriva foi documentada antes da proibição.

Duas características locais importam ao desenho. A primeira é que a cultura
dominante é \textbf{fruta}, não grão: a intensidade de pulverização está ligada a
controle fitossanitário de ciclo curto, com aplicações repetidas ao longo da
safra. A segunda é hidrogeológica: o Aquífero Jandaíra é cárstico, e a
vulnerabilidade à percolação é heterogênea no território --- o que abre variação
explorável no canal-água, e simultaneamente adverte contra tratar todos os
municípios como igualmente permeáveis.

\subsection{O químico, e uma correção que o desenho exige}
\label{sub:quimico}

\begin{quote}
\textbf{Advertência ao leitor familiarizado com a literatura de agrotóxicos e
saúde.} Boa parte dessa literatura --- e, no Brasil, boa parte dos desenhos de
avaliação --- é construída sobre \textbf{glifosato} aplicado a culturas
transgênicas tolerantes a herbicida. \textbf{Não é o caso aqui}, e supor que seja
transferiria mecanismo errado.
\end{quote}

A justificativa do Projeto de Lei nº 18/2015 reproduz tabela de análises de água
oriunda do Dossiê ABRASCO, referente a 23 pontos de coleta na Chapada do Apodi.
A composição encontrada é a seguinte:

\begin{table}[htbp]
\centering
\caption{Princípios ativos detectados em água, Chapada do Apodi}
\label{tab:abrasco}
\begin{tabular}{llc}
\hline
Princípio ativo & Classe & Amostras positivas \\
\hline
Procimidona & Fungicida & 23 de 23 \\
Carbaril & Inseticida (carbamato) & 23 de 23 \\
Carbofurano & Inseticida (carbamato) & 18 de 23 \\
Fenitrotiona & Inseticida (organofosforado) & 16 de 23 \\
Glifosato & Herbicida & 4 de 23 \\
\hline
\end{tabular}

\vspace{0.5em}
\begin{minipage}{0.85\textwidth}
\footnotesize \emph{Fonte:} tabela reproduzida na justificativa do PL nº
18/2015, atribuída ao Dossiê ABRASCO; coleta referente a 2009. \emph{Nota:} a
verificação da fonte primária permanece pendente.
\end{minipage}
\end{table}

O perfil é de \textbf{fungicida e inseticida sobre fruticultura irrigada}, não de
herbicida sobre grão transgênico. Três consequências metodológicas decorrem
disso, e todas são levadas adiante nas seções seguintes:

\begin{enumerate}
  \item \textbf{Os desenhos de referência não transferem quimicamente.} Estudos
    construídos sobre a difusão de soja tolerante a herbicida transferem na
    \emph{forma} --- a lógica de instrumentar exposição por aptidão agroclimática
    ---, não no \emph{mecanismo}. A molécula, a via de exposição e a cultura são
    outras.

  \item \textbf{A receita de instrumentação precisa ser adaptada.} A construção
    do instrumento de aptidão agroclimática, na literatura de referência, toma o
    máximo entre culturas \textbf{transgênicas}. Banana não é transgênica. O
    método transfere; a lista de culturas, não. Verificar se o catálogo
    agroclimático cobre banana e melão é, por isso, pré-requisito e não detalhe.

  \item \textbf{Carbamatos e organofosforados têm assinatura toxicológica
    distinta.} São inibidores de colinesterase, com quadro agudo característico
    --- o que torna a série de internação por intoxicação aguda um canal
    informativo, e não apenas um desfecho secundário.
\end{enumerate}

\subsection{O mecanismo: deriva}
\label{sub:deriva}

O que distingue a aplicação aérea da terrestre é a \textbf{deriva}: parte da
calda não atinge o alvo e é carregada pelo vento, dispersando-se sobre
residências, escolas, corpos d'água e comunidades a favor do vento. A exposição
resultante é \textbf{difusa e populacional}, não ocupacional e pontual --- e é
essa característica que faz da proibição um instrumento de saúde pública, e não
apenas de segurança do trabalho.

Dois canais operam em paralelo, e distingui-los é contribuição de mecanismo:

\begin{description}
  \item[Canal-ar.] Deriva direta sobre populações situadas a favor do vento em
    relação à área de aplicação. A direção do vento é, aqui, fonte potencial de
    identificação: populações a favor e contra o vento diferem em exposição sem
    diferir sistematicamente em outros atributos.

  \item[Canal-água.] Escoamento e percolação para bacias hidrográficas, com
    chegada à água de consumo. A estrutura montante/jusante dentro de bacia
    oferece comparação interna, e a heterogeneidade cárstica do Jandaíra modula a
    intensidade da percolação.
\end{description}

\paragraph{Um terceiro canal, de sinal oposto.} Proibido o avião, a aplicação não
desaparece: migra para trator e pulverizador costal. Isso \emph{afasta} o produto
da população --- menos deriva --- e simultaneamente o \emph{aproxima} do
aplicador. A proibição pode, portanto, reduzir exposição populacional e
\textbf{aumentar} intoxicação aguda ocupacional. O canal de substituição é
medido separadamente, e sua direção esperada é oposta à do desfecho principal.

\subsection{Fiscalização}
\label{sub:enforcement}

A fiscalização do art. 28-B cabe, pelos arts. 15 e 30 da Lei nº 12.228/1993, à
Superintendência Estadual do Meio Ambiente, à Secretaria da Agricultura e à
Secretaria da Saúde. O art. 8º da mesma lei determina o registro, junto ao órgão
ambiental estadual, das empresas prestadoras de serviço que utilizam agrotóxicos
--- registro que, se público e com série histórica municipal, constitui a única
fonte estadual capaz de identificar quem operava aplicação aérea antes da
proibição.

A extensão efetiva da fiscalização não é conhecida. Trata-se de limitação
declarada, com consequência precisa sobre o estimando: se a proibição não foi
fiscalizada, a dose pré-ban mede \emph{intenção} de aplicação, não exposição
efetivamente removida, e o parâmetro estimado passa a ser um efeito de intenção
de tratar. A §\ref{sub:ameacas} trata do ponto, e a §\ref{sub:deriva} indica a
via independente pela qual ele pode ser parcialmente contornado.
